\section{Soluciones existentes de simulación y análisis financiero}

El ecosistema de herramientas relacionadas con inversión y mercados es amplio, pero muy heterogéneo. Bajo la misma etiqueta conviven plataformas de \textit{paper trading}, simuladores bursátiles de orientación formativa, herramientas de análisis de carteras y servicios profesionales centrados en operativa o visualización. Por ello, para situar correctamente la aportación de este TFG no basta con afirmar que “ya existen simuladores”, sino que conviene analizar qué tipo de problema resuelve realmente cada familia de soluciones y qué lagunas dejan abiertas desde el punto de vista pedagógico.

En el caso de este proyecto, el objetivo no era construir un terminal de trading ni una plataforma de inversión real, sino una aplicación web educativa capaz de combinar datos históricos reales, experimentación sin riesgo, progresión guiada, persistencia de usuario e interpretación de resultados. A continuación se revisan cinco referencias representativas que permiten contextualizar esa decisión.

\subsection{Investopedia Simulator}

Investopedia Stock Simulator \cite{investopedia} es probablemente una de las referencias más reconocibles en el terreno del aprendizaje bursátil para usuarios principiantes. Se trata de un simulador basado en dinero virtual que permite crear carteras, ejecutar compras y ventas simuladas y seguir la evolución de posiciones en un entorno inspirado en la operativa de mercado real. Su valor principal reside en reducir la barrera de entrada: el usuario puede practicar sin exponer capital propio y familiarizarse con conceptos como órdenes, posiciones, diversificación o rendimiento acumulado.

Desde un punto de vista funcional, la plataforma ofrece una experiencia cercana al \textit{paper trading}. El usuario dispone de saldo virtual, puede seleccionar activos y observar cómo variaría su cartera si esas operaciones se hubieran realizado en el mercado. Además, el hecho de estar integrada en un ecosistema divulgativo como Investopedia refuerza su valor como puerta de acceso a terminología y conceptos básicos. Para una fase inicial de alfabetización financiera, esta combinación de simulación y material explicativo resulta útil.

Sus ventajas principales son, por tanto, la accesibilidad, la familiaridad de la interfaz y la capacidad de experimentar con decisiones de compra y venta en un entorno relativamente comprensible. También destaca por el valor pedagógico indirecto de trabajar sobre una marca conocida en educación financiera, lo que genera confianza en el usuario principiante.

Sin embargo, sus inconvenientes son relevantes para el enfoque de este TFG. En primer lugar, la lógica dominante sigue siendo transaccional: el usuario aprende sobre todo a operar, no necesariamente a reflexionar sobre una estrategia histórica de largo plazo. En segundo lugar, la experiencia no está diseñada alrededor de una progresión didáctica cerrada con itinerario propio, sino alrededor de la simulación libre de operaciones. En tercer lugar, no incorpora un modo equivalente al \textit{Modo Carrera}, donde el paso del tiempo, los turnos, los eventos y la comparación estructurada con un \textit{benchmark} obligan a reinterpretar decisiones anteriores. Tampoco ofrece una capa de tutoría educativa sobre un informe propio del usuario como la que se desarrolla aquí.

La relación con el presente TFG es clara: ambas soluciones comparten la idea de practicar sin riesgo con datos vinculados al mercado, pero divergen en el tipo de aprendizaje que priorizan. Mientras Investopedia Simulator se acerca más a una práctica libre de operativa y seguimiento de cartera, este proyecto se orienta a una experiencia formativa guiada, más centrada en comprender resultados, contexto histórico y gestión del riesgo que en reproducir la operativa de un bróker.

\subsection{La Bolsa Virtual y simuladores bursátiles en español}

La Bolsa Virtual \cite{bolsavirtual} representa una referencia relevante dentro del contexto hispanohablante. Este tipo de plataformas ofrece al usuario la posibilidad de simular inversiones sobre mercados reales, con compra y venta de activos, seguimiento de carteras y, en algunos casos, componentes de comunidad o competición entre participantes. Desde una perspectiva académica, resultan interesantes porque trasladan al ámbito local un modelo de aprendizaje por experimentación que no depende exclusivamente de herramientas anglosajonas.

Lo que ofrecen estas soluciones es, ante todo, un entorno de simulación práctica con terminología, mercados y estilo de uso cercanos al pequeño inversor minorista. Para una persona que quiera entender cómo evoluciona una cartera o qué significa operar con determinados valores, este tipo de simuladores puede servir como punto de partida razonable. Su principal fortaleza es la cercanía de idioma y contexto, algo especialmente valioso cuando el objetivo es introducir conceptos a estudiantes o usuarios sin experiencia previa.

Entre sus ventajas se encuentran precisamente esa accesibilidad lingüística, el enfoque de aprendizaje sin riesgo real y la posibilidad de acostumbrarse a dinámicas de mercado sin abrir una cuenta de inversión auténtica. También pueden ayudar a visualizar de forma intuitiva el impacto de las decisiones de compra y venta sobre el valor de una cartera.

No obstante, sus limitaciones son importantes cuando se comparan con los objetivos de este TFG. Muchas de estas plataformas siguen centradas en reproducir el marco mental del mercado y del bróker, no en construir un recorrido pedagógico progresivo. Además, el foco suele situarse en el seguimiento de operaciones y resultados, más que en el análisis razonado del proceso de decisión. La combinación entre teoría, validación de conocimientos, simulación puntual, simulación por turnos y lectura guiada del informe final no suele aparecer integrada en un mismo producto.

En ese sentido, la relación con este proyecto es de afinidad parcial. Comparten la idea de practicar inversión con datos vinculados al mercado y sin exposición económica real, pero no cubren de forma equivalente la dimensión de progresión educativa, diferenciación entre usuario autenticado e invitado, persistencia del aprendizaje o tutoría interpretativa opcional sobre resultados. La propuesta desarrollada en el TFG intenta ir un paso más allá al articular un ecosistema formativo completo, no solo un simulador aislado.

\subsection{Paper trading de brókers y TradingView}

Otra familia de referencia está compuesta por los entornos de \textit{paper trading} integrados en plataformas de mercado o análisis, como el \textit{paper trading} de TradingView \cite{tradingview_paper} o las cuentas demo de determinados brókers. En estos casos, la lógica es aún más cercana a la operativa real: el usuario observa gráficos, define entradas y salidas, gestiona posiciones y reproduce de forma casi inmediata decisiones similares a las que tomaría en una cuenta de inversión auténtica.

Estas herramientas ofrecen una experiencia muy valiosa para usuarios interesados en familiarizarse con plataformas de análisis técnico, ejecución simulada y visualización avanzada del mercado. Su principal fortaleza es la cercanía al mundo profesional o semiprofesional del \textit{trading}. Permiten entender órdenes, tiempos de ejecución, seguimiento intradía o construcción visual de hipótesis operativas a partir de gráficos.

Precisamente por ello, sus ventajas no coinciden necesariamente con las prioridades de este TFG. Son herramientas potentes, maduras y útiles cuando el objetivo es aprender a usar interfaces de análisis de mercado o practicar operativa simulada. Además, plataformas como TradingView destacan por la calidad de su visualización, la amplitud de instrumentos y la comunidad de usuarios.

Sin embargo, desde el punto de vista de un trabajo académico orientado a educación financiera de largo plazo, presentan varios inconvenientes. El primero es que su centro de gravedad conceptual suele estar más próximo al análisis gráfico, al seguimiento del precio y a la operativa que a la comprensión pedagógica del proceso inversor completo. El segundo es que no suelen organizar la experiencia en modos diferenciados como práctica puntual, progreso previo, carrera histórica o proyección experimental. El tercero es que no integran de forma natural un informe final didáctico propio del sistema ni una capa de interpretación educativa acotada como el Tutor IA desarrollado aquí.

Su relación con este proyecto es importante porque representan el polo contrario del diseño perseguido. Mientras el \textit{paper trading} clásico intenta parecerse a un entorno real de mercado, el TFG busca deliberadamente una experiencia más guiada, más accesible y menos orientada a la operativa rápida. La diferencia no implica que una solución sea mejor que la otra en términos absolutos, sino que responden a problemas distintos. En este caso, la prioridad académica era reforzar comprensión, contexto y reflexión, no reproducir con la máxima fidelidad una mesa de negociación o una plataforma de \textit{trading} técnico.

\subsection{Portfolio Visualizer y herramientas de backtesting de carteras}

Portfolio Visualizer \cite{portfolio_visualizer} representa otra categoría distinta: la de herramientas centradas en análisis cuantitativo, backtesting y comparación de carteras. En lugar de poner el énfasis en ejecutar órdenes o seguir una cuenta simulada, estas soluciones permiten estudiar comportamientos históricos, comparar asignaciones, medir rentabilidades, volatilidad, drawdown o sensibilidad de una cartera ante distintos escenarios temporales.

Este tipo de herramientas ofrece un valor claro para usuarios que desean analizar decisiones de asignación con mayor profundidad cuantitativa. Sus ventajas son relevantes: permiten trabajar con métricas objetivas, comparar estrategias y obtener una visión más estructurada del comportamiento histórico de una cartera o combinación de activos. En ese sentido, se acercan más al espíritu analítico del presente proyecto que los simuladores de operativa pura.

No obstante, también presentan inconvenientes desde un enfoque pedagógico generalista. A menudo exigen un nivel de alfabetización financiera y estadística mayor, lo que puede convertir la herramienta en un entorno menos accesible para usuarios novatos. Además, suelen centrarse en el análisis del resultado, no tanto en la experiencia narrativa de tomar decisiones sucesivas dentro de un recorrido guiado. Tampoco acostumbran a integrar capas de gamificación, persistencia de progresión educativa o interpretación tutorial sobre el informe final.

La relación con el TFG es doble. Por un lado, constituyen una referencia clara para justificar la importancia de métricas como rentabilidad, \textit{benchmark}, drawdown o tracking error, ya que demuestran que el análisis histórico comparativo es una práctica consolidada. Por otro lado, muestran precisamente la brecha que esta memoria intenta cubrir: traducir parte de ese valor analítico a una experiencia más pedagógica, más accesible y más integrada en una aplicación web orientada a aprendizaje progresivo.

\subsection{Herramientas profesionales de datos y análisis financiero}

Finalmente, conviene considerar el horizonte de herramientas profesionales, como Bloomberg Terminal, Refinitiv Workspace o servicios de datos financieros de gama alta \cite{bloomberg,refinitiv}. Estas soluciones no son comparables al TFG en escala ni en propósito, pero sí resultan útiles como referencia extrema para entender qué significa trabajar con un entorno maduro de información financiera.

Lo que ofrecen estas plataformas es acceso profundo a datos, noticias, métricas, series temporales, análisis, cribado y monitorización de mercado. Sus ventajas son evidentes: amplitud de cobertura, robustez, fiabilidad relativa, herramientas analíticas avanzadas y presencia consolidada en ámbitos profesionales. Desde un punto de vista técnico, representan un estándar alto en integración de información y capacidad de análisis.

Sin embargo, precisamente por su complejidad, coste y orientación profesional, quedan fuera del alcance razonable de un TFG de estas características. No son soluciones concebidas como herramienta introductoria de aprendizaje para estudiantes ni como entorno de simulación gamificada. Su incorporación directa habría sido inviable tanto por coste como por sobredimensión del problema. Además, su uso no encaja con un producto web ligero desplegado en un entorno académico con presupuesto muy limitado.

La comparación con este proyecto sirve para delimitar expectativas. El objetivo del TFG no era competir con proveedores profesionales, sino construir una herramienta educativa funcional que pudiera desplegarse con medios realistas, apoyándose en datos históricos accesibles y en una arquitectura razonablemente mantenible. En ese marco, la elección de tecnologías y fuentes de datos debía responder al equilibrio entre viabilidad, valor pedagógico y trazabilidad técnica.

\subsection{Comparativa y diferenciación de la propuesta propia}

La Tabla~\ref{tab:comparativa_soluciones_existentes} resume la comparación entre las soluciones revisadas y la propuesta desarrollada en este TFG.

\begin{table}[htbp]
\centering
\small
\renewcommand{\arraystretch}{1.15}
\begin{tabular}{|p{0.14\textwidth}|p{0.08\textwidth}|p{0.11\textwidth}|p{0.10\textwidth}|p{0.09\textwidth}|p{0.12\textwidth}|p{0.09\textwidth}|p{0.19\textwidth}|}
\hline
\textbf{Herramienta} & \textbf{Datos reales} & \textbf{Simulación largo plazo} & \textbf{Enfoque educativo} & \textbf{Persistencia} & \textbf{Gamificación / carrera} & \textbf{IA educativa} & \textbf{Limitaciones principales} \\ \hline
Investopedia Simulator & Sí & Parcial & Medio & Sí & No & No & Enfoque más transaccional que pedagógico progresivo. \\ \hline
La Bolsa Virtual & Sí & Parcial & Medio & Sí & Parcial & No & Más cercano a simulador bursátil clásico que a itinerario formativo integrado. \\ \hline
Paper trading de brókers / TradingView & Sí & Baja & Baja & Sí & No & No & Orientado a operativa y análisis gráfico, no a progresión educativa. \\ \hline
Portfolio Visualizer & Sí & Alta & Media & Limitada & No & No & Muy analítico, pero menos guiado y menos accesible para usuarios novatos. \\ \hline
Herramientas profesionales & Sí & Alta & Baja & Sí & No & No & Coste, complejidad y sobredimensión para un TFG. \\ \hline
Proyecto desarrollado & Sí & Alta & Alta & Sí & Sí & Sí, opcional & Dependencia de Yahoo Finance, alcance académico y validación IA principalmente cualitativa. \\ \hline
\end{tabular}
\caption{Comparativa entre soluciones existentes y la propuesta desarrollada en el TFG.}
\label{tab:comparativa_soluciones_existentes}
\end{table}

A la vista de esta comparación, la principal diferenciación del proyecto no reside en inventar un tipo de herramienta completamente nuevo, sino en integrar dentro de una misma aplicación varios elementos que normalmente aparecen separados: aprendizaje progresivo, simulación histórica, persistencia por usuario, modo carrera gamificado, comparación con benchmark, proyección experimental y una capa opcional de interpretación educativa apoyada por IA. Esa combinación es precisamente la que justifica la propuesta como desarrollo propio y como objeto de análisis académico diferenciado.

\section{Alternativas tecnológicas evaluadas}

La elección tecnológica del proyecto no fue arbitraria. Aunque el sistema final utiliza Flask, PostgreSQL, Peewee, Jinja2, Chart.js, yfinance, Render y GitHub Actions, esa combinación se entiende mejor cuando se compara con otras alternativas razonables que podrían haberse adoptado. El interés de este apartado no es reconstruir una falsa fase de selección exhaustiva inexistente, sino justificar por qué las decisiones finales resultan coherentes con el alcance, el presupuesto y la evolución real del TFG.

\subsection{Backend web: Flask frente a Django, FastAPI y Streamlit}

En el plano del backend, una primera decisión fundamental era escoger el marco principal de la aplicación. \textbf{Flask} \cite{flask} ofrece un enfoque de microframework, con muy pocas imposiciones estructurales y una curva de entrada suave para proyectos web que necesitan combinar rutas HTML, endpoints JSON y lógica de negocio propia. Esa flexibilidad encaja bien con un TFG que debía evolucionar por iteraciones y que no partía de un dominio corporativo masivo ni de un modelo administrativo complejo.

Frente a Flask, \textbf{Django} \cite{django} representa una solución más completa y opinionada. Sus ventajas son claras: panel administrativo integrado, ORM potente, sistema de autenticación maduro y una gran cantidad de componentes listos para usar. Para aplicaciones de negocio con más entidades, flujos de backoffice o necesidades de administración avanzadas, Django habría sido una opción muy defendible. Sin embargo, también introduce mayor peso estructural, más convenciones y una inercia de proyecto menos ligera. Para el alcance real del TFG, esa sobrecarga habría incrementado el coste cognitivo y de integración.

\textbf{FastAPI} \cite{fastapi} habría sido una alternativa interesante si el objetivo principal hubiera sido construir una API moderna, fuertemente tipada y orientada a consumo por frontend desacoplado. Entre sus ventajas destacan el tipado, la generación automática de documentación y el buen rendimiento en servicios HTTP. No obstante, la aplicación entregada no necesitaba una API pública compleja ni un frontend separado de forma estricta. Gran parte del valor del proyecto reside precisamente en el renderizado servidor con plantillas y en una integración progresiva de JavaScript, por lo que FastAPI no ofrecía una ventaja decisiva para este caso.

Por último, \textbf{Streamlit} \cite{streamlit} podría haber permitido construir con rapidez una interfaz orientada a datos y visualización, muy útil en prototipos de análisis o demostraciones. Su principal ventaja es la velocidad de prototipado para herramientas de ciencia de datos. Sin embargo, sus inconvenientes para este proyecto eran importantes: menor control fino sobre la arquitectura web, menor adecuación para autenticación personalizada, dificultad relativa para sostener una experiencia multipantalla con navegación y menor encaje con un producto web que debía comportarse como una aplicación más generalista.

La elección de Flask se justifica, por tanto, por equilibrio. Permite servir HTML, exponer JSON, integrar autenticación propia, mantener el proyecto razonablemente simple y sostener una arquitectura monolítica ligera adecuada para el TFG. No ofrece tantas piezas integradas como Django ni tantas garantías tipadas como FastAPI, pero precisamente esa moderación encaja mejor con el problema resuelto.

\subsection{Persistencia: SQLite, PostgreSQL, MySQL/MariaDB y almacenamiento JSON}

La persistencia era otro punto decisivo, especialmente cuando el proyecto dejó de ser un prototipo puntual y empezó a necesitar autenticación, historial y sesiones recuperables. \textbf{SQLite} \cite{sqlite} suele ser una elección razonable en fases tempranas por su sencillez operativa: no requiere servicio externo, reduce fricción local y es muy útil para pruebas o desarrollos personales. Su principal ventaja es precisamente esa ligereza. Sin embargo, en un despliegue real con hosting externo y necesidad de continuidad entre reinicios, sus límites aparecen con rapidez. En este proyecto se aceptó SQLite como base efímera de integración continua, pero no como solución final de producto.

\textbf{PostgreSQL} \cite{postgresql} se eligió como solución principal de persistencia en Render. Sus ventajas son conocidas: robustez, madurez, buen soporte en la nube y un modelo relacional suficientemente sólido para usuarios, historial y sesiones del modo carrera. Para el alcance del TFG no era necesario explotar características avanzadas del motor, pero sí disponer de una base fiable y compatible con un despliegue real. Además, Render ofrece integración razonable con este tipo de base de datos, lo que reducía la complejidad de despliegue.

\textbf{MySQL/MariaDB} \cite{mariadb} también habría sido una alternativa válida. Sus ventajas incluyen madurez, amplia implantación y abundante documentación. Desde el punto de vista del TFG, no existía una razón funcional fuerte para descartarlas por incapacidad técnica, pero sí había menos alineación directa con la solución que finalmente se estabilizó en Render y con la configuración ya empleada en el proyecto. En otras palabras, habrían sido viables, pero no claramente mejores para este caso.

Por último, el proyecto utiliza también almacenamiento \textbf{JSON/local} en ciertos puntos concretos, como el store operativo del modo carrera o el ranking local. Esta opción tiene una ventaja pragmática: permite conservar compatibilidad incremental y simplificar parte de la operativa sin rediseñar por completo el motor existente. Su inconveniente es igualmente claro: no debe confundirse con la capa persistente principal ni con una solución relacional robusta para producción.

La elección final de PostgreSQL en Render se justifica porque era la opción que mejor equilibraba realismo, fiabilidad y viabilidad académica. El uso de JSON quedó relegado a apoyo operativo parcial y a compatibilidad incremental, mientras que SQLite se mantuvo como herramienta válida de CI, no como base final de producto.

\subsection{Capa de acceso a datos: Peewee frente a SQLAlchemy y SQL directo}

Una vez elegida la persistencia relacional, era necesario definir cómo modelar el acceso a datos. El proyecto final emplea \textbf{Peewee} \cite{peewee}, un ORM ligero para Python. Su principal ventaja en este contexto es la simplicidad: permite definir modelos claros, relaciones básicas y operaciones frecuentes sin introducir una capa excesivamente compleja. Para una aplicación académica de tamaño contenido, Peewee resulta suficientemente expresivo y más fácil de dominar que soluciones más pesadas.

La alternativa más evidente habría sido \textbf{SQLAlchemy} \cite{sqlalchemy}, probablemente el ORM más extendido en proyectos Python de cierta entidad. Sus ventajas son poderosas: gran flexibilidad, ecosistema amplio, madurez, expresividad y capacidad para crecer con arquitecturas más complejas. Sin embargo, también implica mayor complejidad conceptual y más decisiones de configuración. En el contexto de este TFG, esa potencia adicional no era imprescindible para resolver los modelos efectivamente usados.

Otra posibilidad habría sido recurrir a \textbf{SQL directo} con consultas manuales. Esta opción tiene una ventaja teórica importante: control absoluto sobre las consultas y ausencia de abstracciones intermedias. No obstante, sus inconvenientes para un proyecto como este son claros: más código repetitivo, mayor propensión a errores, más esfuerzo de mantenimiento y peor expresividad a la hora de representar entidades como usuario, historial, sesiones o resultados del readiness quiz.

La elección de Peewee se justifica así por proporción. Ofrece una abstracción suficiente para modelar \texttt{User}, \texttt{AnalysisHistory}, \texttt{CareerSession}, \texttt{CareerTurn}, \texttt{CareerSessionLink} y \texttt{ReadinessQuizResult} sin sobredimensionar el diseño. Además, esta memoria debe dejar expresamente claro que el proyecto \textbf{no usa SQLAlchemy}, ya que una de las revisiones detectó el riesgo de describir el stack con etiquetas genéricas no fieles al repositorio real.

\subsection{Frontend: Jinja2 y JavaScript progresivo frente a SPA completa}

En el plano de la interfaz, el proyecto adopta un enfoque basado en \textbf{Jinja2} \cite{jinja2} para el renderizado en servidor y \textbf{JavaScript progresivo} para enriquecer interacciones concretas. Esta decisión tiene varias ventajas: simplifica el despliegue, reduce el número de piezas tecnológicas, facilita la consistencia entre sesión y renderizado HTML y evita el coste de construir una SPA completa solo para resolver una aplicación cuyo núcleo sigue descansando en pantallas de servidor.

Alternativas como \textbf{React} o \textbf{Vue} \cite{react,vue} habrían permitido una experiencia de cliente más desacoplada y probablemente una gestión más sofisticada del estado de interfaz. Sus ventajas son bien conocidas: ecosistema amplio, componentes reutilizables y una gran capacidad para construir aplicaciones muy dinámicas. Sin embargo, también habrían exigido más infraestructura de desarrollo, más complejidad de build y una separación frontend/backend que no era indispensable para el alcance académico del proyecto.

Una \textbf{SPA completa} también habría sido defendible si el problema principal hubiera sido la interacción continua del cliente con una API rica y fuertemente desacoplada. Pero el TFG combina páginas informativas, formularios, rutas renderizadas, autenticación sencilla y llamadas JSON puntuales. En ese marco, una SPA completa habría introducido un coste estructural difícil de justificar frente al valor adicional real.

En cuanto a la visualización, se optó por \textbf{Chart.js} \cite{chartjs} frente a alternativas como \textbf{D3.js} \cite{d3}. Chart.js ofrece rapidez de integración, buena relación entre simplicidad y expresividad y suficiente capacidad para las gráficas de análisis, benchmark y series del modo carrera. D3.js, en cambio, proporciona un control mucho más fino y poderoso, pero a costa de una complejidad mayor. Para una aplicación cuyo objetivo no era innovar en visualización sino presentar datos de forma clara y mantenible, Chart.js resultaba más proporcionado.

La elección final de Jinja2 más JavaScript progresivo y Chart.js se justifica por coherencia con una arquitectura ligera, más fácil de mantener, plenamente compatible con Flask y suficiente para el nivel de interacción realmente necesario.

\subsection{Fuentes de datos financieros: Yahoo Finance/yfinance frente a Alpha Vantage, FMP y proveedores profesionales}

La obtención de datos financieros reales era un requisito central. El proyecto emplea \textbf{Yahoo Finance} a través de \textbf{yfinance} \cite{yfinance}. Su ventaja principal es muy clara: acceso relativamente sencillo y sin coste directo obligatorio a históricos ajustados suficientes para un proyecto académico. Esta disponibilidad permite construir backtests, métricas y comparativas sin introducir desde el inicio un coste económico alto ni barreras de integración mayores.

Alternativas como \textbf{Alpha Vantage} \cite{alphavantage} o \textbf{Financial Modeling Prep} \cite{fmp} ofrecen APIs documentadas y, en algunos casos, modalidades gratuitas o freemium. Sus ventajas incluyen endpoints más formalizados y una relación API-consumidor más explícita. Sin embargo, también suelen implicar límites de uso, claves específicas, dependencia de planes comerciales para crecer y un coste potencialmente menos conveniente si el objetivo es sostener muchos experimentos o iteraciones en un contexto académico limitado.

En el otro extremo, proveedores profesionales como Bloomberg, Refinitiv o Nasdaq Data Link \cite{bloomberg,refinitiv,nasdaq_data_link} proporcionan más robustez, cobertura y soporte. Pero su coste y complejidad los hacen desproporcionados para el TFG. Además, el problema del proyecto no exigía una terminal profesional, sino una fuente suficientemente razonable para simulación histórica educativa.

La elección de yfinance se justifica, por tanto, por coste, accesibilidad e integración. No obstante, esta memoria debe reconocer sus límites: dependencia de un proveedor externo no oficial, posibilidad de respuestas vacías, rate limiting, cambios no controlados y necesidad de degradación amistosa ante fallos. En otras palabras, la elección es adecuada para el alcance académico, pero no debe presentarse como equivalente a una fuente profesional garantizada.

\subsection{Despliegue e integración continua: Render y GitHub Actions frente a otras opciones}

El despliegue final del sistema se estabilizó en \textbf{Render} \cite{render}, mientras que la validación automatizada se apoyó en \textbf{GitHub Actions} \cite{github_actions}. Esta combinación ofrece varias ventajas para un TFG: integración relativamente directa con GitHub, facilidad de despliegue de una aplicación Flask con Gunicorn, soporte razonable para variables de entorno y una barrera de entrada mucho menor que la de proveedores de infraestructura más complejos.

Alternativas como \textbf{Railway} \cite{railway}, los antiguos flujos de \textbf{Heroku} \cite{heroku}, o despliegues más manuales sobre \textbf{AWS} \cite{aws} habrían sido posibles. Cada una tiene sus ventajas. Railway destaca por sencillez y rapidez de puesta en marcha; AWS ofrece escalabilidad y un ecosistema enorme; incluso una validación puramente local sin CI/CD reduce complejidad inicial. Sin embargo, para el objetivo del proyecto importaba más disponer de una ruta de despliegue realista y defendible que de una infraestructura sobredimensionada.

GitHub Actions, por su parte, aporta un mínimo de disciplina de calidad: estilo, formato y pruebas automatizadas sobre el repositorio. Podría haberse prescindido de CI y limitarse a validación local, pero eso habría debilitado la trazabilidad técnica del cierre. Incorporar Actions no convierte el proyecto en una plataforma industrial, pero sí aporta una señal clara de madurez metodológica.

La elección final de Render más GitHub Actions se justifica por equilibrio entre sencillez, coste y valor demostrable. No es la solución más potente imaginable, pero sí una de las más adecuadas para entregar una aplicación web académica desplegada, con base de datos, validación automatizada y flujo de trabajo suficientemente profesional para el alcance perseguido.

\subsection{Síntesis de la elección tecnológica}

En conjunto, la arquitectura final del proyecto responde a una lógica de proporcionalidad. Flask permite una base web ligera y flexible; PostgreSQL aporta persistencia real; Peewee ofrece una capa de datos suficiente sin sobredimensionar el ORM; Jinja2 y Chart.js resuelven la interfaz y la visualización con complejidad moderada; yfinance proporciona históricos accesibles; y Render junto con GitHub Actions permiten despliegue y validación defendibles.

La elección no debe interpretarse como la única combinación técnicamente válida, sino como la más coherente con un TFG que debía llegar a un producto funcional, desplegado y explicable, sin perder de vista el tiempo, el presupuesto y la necesidad de mantener la complejidad bajo control.
